% ============================================================
% Beamer Header — matches stats-report-template.tex palette
% K. El Husseini
% ============================================================

% ----- TikZ -----
\usepackage{tikz}
\usetikzlibrary{positioning,calc,arrows.meta,shapes.geometric}

% ----- Code blocks : taille de police -----
% fancyvrb est chargé conditionnellement par Pandoc (seulement si code blocks présents).
% On le charge ici pour garantir que \DefineVerbatimEnvironment est toujours disponible.
\usepackage{fancyvrb}
\DefineVerbatimEnvironment{Highlighting}{Verbatim}{fontsize=\footnotesize,commandchars=\\\{\}}

% ----- Theme base -----
\usetheme{default}
\useinnertheme{rectangles}
\usefonttheme{professionalfonts}  % prevent Beamer overriding fontspec fonts

% ----- Fonts -----
\usepackage{fontspec}
\setmainfont{Inter}[
  UprightFont     = *-Regular,
  BoldFont        = *-Bold,
  ItalicFont      = *-Italic,
  BoldItalicFont  = *-BoldItalic,
  Extension       = .otf,
  Path            = /Users/kinelhu/Library/Fonts/
]
\setsansfont{Inter}[
  UprightFont     = *-Regular,
  BoldFont        = *-Bold,
  ItalicFont      = *-Italic,
  Extension       = .otf,
  Path            = /Users/kinelhu/Library/Fonts/
]
\setmonofont{JetBrains Mono}[Scale=0.85]

% ----- Colors -----
% Seeds : seules ces 3 valeurs sont à modifier pour changer le thème.
% Tout le reste est dérivé automatiquement via xcolor.
\definecolor{AccentPrimary}{HTML}{27A9BC}  % teinte principale default: 0D7377
\definecolor{AlertOrange}{HTML}{284BBD}    % avertissement / alertblock default: C0550A
\definecolor{ExampleGreen}{HTML}{28BD9F}   % exemple positif / exampleblock default: 1A6B3C

% Dérivées de AccentPrimary
\colorlet{AccentDark}{AccentPrimary!60!black}   % variante sombre
\colorlet{AccentLight}{AccentPrimary!7!white}   % tinte très claire (titre page)
\colorlet{CalloutBg}{AccentPrimary!6!white}     % fond des block standard

% Dérivées de AlertOrange et ExampleGreen
\colorlet{AlertOrangeBg}{AlertOrange!8!white}
\colorlet{ExampleGreenBg}{ExampleGreen!8!white}

% Neutres (indépendants de la teinte)
\colorlet{TextPrimary}{black!90}
\colorlet{TextSecondary}{black!42}
\colorlet{BorderLight}{black!12}
\colorlet{PageBackground}{white}

% ----- Beamer colors -----
\setbeamercolor{normal text}{fg=TextPrimary, bg=PageBackground}
\setbeamercolor{structure}{fg=AccentPrimary}
% alerted text / example text : pilotent local structure à l'intérieur des blocs.
% Beamer fait \setbeamercolor{local structure}{parent=alerted text} en entrant dans
% un alertblock — il faut donc que alerted text porte la bonne couleur.
\setbeamercolor{alerted text}{fg=AlertOrange}
\setbeamercolor{example text}{fg=ExampleGreen}
\setbeamercolor{title}{fg=white, bg=AccentPrimary}
\setbeamercolor{subtitle}{fg=AccentLight, bg=AccentPrimary}
\setbeamercolor{author}{fg=AccentLight, bg=AccentPrimary}
\setbeamercolor{date}{fg=AccentLight, bg=AccentPrimary}
\setbeamercolor{frametitle}{fg=AccentPrimary, bg=PageBackground}
\setbeamercolor{block title}{fg=white, bg=AccentPrimary}
\setbeamercolor{block body}{fg=TextPrimary, bg=CalloutBg}
\setbeamercolor{block title alerted}{fg=white, bg=AlertOrange}
\setbeamercolor{block body alerted}{fg=TextPrimary, bg=AlertOrangeBg}
\setbeamercolor{block title example}{fg=white, bg=ExampleGreen}
\setbeamercolor{block body example}{fg=TextPrimary, bg=ExampleGreenBg}
\setbeamercolor{itemize item}{fg=AccentPrimary}
\setbeamercolor{itemize subitem}{fg=TextSecondary}
\setbeamercolor{enumerate item}{fg=AccentPrimary}
\setbeamercolor{footline}{fg=TextSecondary, bg=BorderLight}

% ----- Beamer fonts -----
\setbeamerfont{title}{size=\LARGE, series=\bfseries}
\setbeamerfont{subtitle}{size=\large, series=\mdseries}
\setbeamerfont{frametitle}{size=\Large, series=\bfseries}
\setbeamerfont{author}{size=\normalsize}
\setbeamerfont{date}{size=\small}
\setbeamerfont{footline}{size=\tiny}
\setbeamerfont{block title}{size=\normalsize, series=\bfseries}
\setbeamerfont{section in toc}{size=\large, series=\bfseries}

% ----- Per-slide logo state (déclaré avant le template frametitle qui s'en sert) -----
\def\currentslidelogo{}
\def\currentslidelogoheight{1.5cm}

% ----- Frame title: colored text + thin underline -----
% \smash[b] zeroes the depth of the title box so descenders (j, p, g...)
% don't push the rule down — consistent baseline-to-rule distance always.
% Logo optionnel : aligné à droite sur la ligne du titre via \hfill.
\setbeamertemplate{frametitle}{%
  \vspace{0.4em}%
  \usebeamerfont{frametitle}\usebeamercolor[fg]{frametitle}%
  \smash[b]{\insertframetitle}%
  \ifx\currentslidelogo\empty\else
    \hfill\smash[b]{\RawIncludegraphics[height=\currentslidelogoheight]{\currentslidelogo}}%
    \global\let\currentslidelogo\empty%
  \fi
  \par%
  \vspace{0.15em}%
  {\color{AccentPrimary}\hrule height 0.4pt}%
  \vspace{0.2em}%
}

% ----- Section divider slides: badge style -----
\usepackage{tcolorbox}
\tcbuselibrary{skins}
\AtBeginSection[]{%
  \begin{frame}[plain]
    \vfill
    \centering
    \begin{tcolorbox}[
      enhanced,
      colback=AccentPrimary,
      colframe=AccentPrimary,
      arc=4pt,
      boxrule=0pt,
      left=16pt, right=16pt, top=6pt, bottom=6pt,
      halign=center,
      width=0.85\textwidth
    ]
      {\usebeamerfont{title}\color{white}\insertsectionhead\par}
    \end{tcolorbox}
    \vfill
  \end{frame}
}

% ----- Navigation off -----
\setbeamertemplate{navigation symbols}{}

% ----- Bullet styling -----
% itemize item : carré plein coloré via local structure.
% local structure hérite de alerted text / example text à l'intérieur des blocs
% (mécanisme natif Beamer — cf. beamerbaselocalstructure.sty).
% Pas de hooks \addtobeamertemplate nécessaires : la chaîne d'héritage suffit.
\setbeamertemplate{itemize item}{%
  \usebeamercolor[fg]{local structure}\rule[0.3ex]{0.5em}{0.5em}}
\setbeamertemplate{itemize subitem}{%
  \usebeamercolor[fg]{itemize subitem}\small\textbullet}
\setbeamertemplate{itemize subsubitem}{%
  \usebeamercolor[fg]{itemize subitem}\tiny\textbullet}
% item projected (carré des listes numérotées) : bg=item.fg→local structure.fg
% → AccentPrimary hors blocs, AlertOrange / ExampleGreen dans les blocs.

% ----- Footer -----
% Trois boîtes égales : auteur (gauche) | date (centré) | n°/total (droite)
% \hfill seul ne centre pas si les côtés sont de longueurs inégales.
\newlength{\beamerprogresswidth}
\setbeamertemplate{footline}{%
  \leavevmode%
  \hbox{%
    \begin{beamercolorbox}[wd=.40\paperwidth, ht=2.4ex, dp=1.2ex,
        left, leftskip=1em]{footline}%
      \usebeamerfont{footline}\insertshortauthor%
    \end{beamercolorbox}%
    \begin{beamercolorbox}[wd=.20\paperwidth, ht=2.4ex, dp=1.2ex,
        center]{footline}%
      \usebeamerfont{footline}\insertshortdate%
    \end{beamercolorbox}%
    \begin{beamercolorbox}[wd=.40\paperwidth, ht=2.4ex, dp=1.2ex,
        right, rightskip=1em]{footline}%
      \usebeamerfont{footline}\insertframenumber\,/\,\inserttotalframenumber%
    \end{beamercolorbox}%
  }%
  \ifnum\inserttotalframenumber>1
    \setlength{\beamerprogresswidth}{\paperwidth}%
    \divide\beamerprogresswidth by \inserttotalframenumber%
    \multiply\beamerprogresswidth by \insertframenumber%
    \begin{beamercolorbox}[wd=\paperwidth, ht=0.4pt]{structure}%
      \textcolor{AccentPrimary}{\rule{\beamerprogresswidth}{0.4pt}}%
    \end{beamercolorbox}%
  \fi
}

% ----- Title page -----
\setbeamertemplate{title page}{%
  \vspace*{2.5em}%
  \begin{center}
    {\usebeamerfont{title}\color{AccentPrimary}\inserttitle\par}%
    \vspace{0.5em}
    {\color{AccentPrimary}\rule{0.65\textwidth}{0.8pt}\par}%
    \vspace{0.5em}
    \ifx\insertsubtitle\@empty\else
      {\usebeamerfont{subtitle}\color{TextSecondary}\insertsubtitle\par}%
    \fi
  \end{center}
  \vfill
  \begin{center}
    {\usebeamerfont{author}\color{TextPrimary}\insertauthor\par}%
    \vspace{0.3em}
    {\usebeamerfont{date}\color{TextSecondary}\insertdate\par}%
  \end{center}
  \vspace{1.5em}%
}

% ----- Caption styling -----
\setbeamertemplate{caption}{\centering\tiny\color{TextSecondary}\insertcaption\par}

% ----- Image handling (pandoc 3.x) -----
\usepackage{adjustbox}
% Auto-center explicit-width images; flag prevents double-centering inside \pandocbounded
\newif\ifinpandocbounded
\let\RawIncludegraphics\includegraphics
\renewcommand{\includegraphics}[2][]{%
  \ifinpandocbounded
    \RawIncludegraphics[#1]{#2}%
  \else
    \begin{center}\RawIncludegraphics[#1]{#2}\end{center}%
  \fi
}
\def\pandocbounded#1{%
  \inpandocboundedtrue
  {\centering\adjustbox{max width=0.85\linewidth, max totalheight=0.65\textheight}{#1}\par}%
  \inpandocboundedfalse
}

% ----- List font size shorthands (use inside \begingroup...\endgroup) -----
% Usage: \begingroup \fnsize\n\n- list\n\n\endgroup
\newcommand{\smsize}{\setbeamerfont{itemize/enumerate body}{size*={9}{11}}}
\newcommand{\fnsize}{\setbeamerfont{itemize/enumerate body}{size*={8}{10}}}
\newcommand{\sssize}{\setbeamerfont{itemize/enumerate body}{size*={7}{9}}}

% ----- Font size environments (for ::: {.smtext} etc.) -----
% Uses setbeamerfont with size*={pt}{leading} so Beamer item templates
% (coloured boxes, bullets) scale correctly alongside the text.
% Changes are group-scoped via \begin/\end and restore automatically.
\newenvironment{smtext}{%
  \setbeamerfont{normal text}{size*={9}{11}}%
  \setbeamerfont{itemize item}{size*={9}{11}}%
  \setbeamerfont{enumerate item}{size*={9}{11}}%
  \setbeamerfont{item projected}{size*={9}{11}}%
  \usebeamerfont{normal text}%
}{}
\newenvironment{fntext}{%
  \setbeamerfont{normal text}{size*={8}{10}}%
  \setbeamerfont{itemize item}{size*={8}{10}}%
  \setbeamerfont{enumerate item}{size*={8}{10}}%
  \setbeamerfont{item projected}{size*={8}{10}}%
  \usebeamerfont{normal text}%
}{}
\newenvironment{sstext}{%
  \setbeamerfont{normal text}{size*={7}{9}}%
  \setbeamerfont{itemize item}{size*={7}{9}}%
  \setbeamerfont{enumerate item}{size*={7}{9}}%
  \setbeamerfont{item projected}{size*={7}{9}}%
  \usebeamerfont{normal text}%
}{}

% ----- Pandoc compat -----
\providecommand{\tightlist}{%
  \setlength{\itemsep}{0pt}\setlength{\parskip}{0pt}}

% ----- Per-slide logo (dans la zone de titre, aligné à droite) -----
% \slidelogo{img/logo.png}      → logo hauteur 1.5 cm (défaut)
% \slidelogo[2cm]{img/logo.png} → hauteur personnalisée
% Appeler dans le corps du frame (pas dans le titre ##).
% Le rendu se fait dans \setbeamertemplate{frametitle} via \hfill — aucun package requis.
% Réinitialisation automatique frame par frame (dans le template frametitle).
\newcommand{\slidelogo}[2][1.5cm]{%
  \global\def\currentslidelogoheight{#1}%
  \global\def\currentslidelogo{#2}%
}
